La varietà di dispositivi e paradigmi disponibili per la raccolta di informazioni cerebrali e la rapidità nello sviluppo di nuovi algoritmi sempre più complessi rendono essenziale il confronto continuo con la letteratura attuale. Consapevole di ciò, la comunità scientifica rinnova costantemente le \emph{review} \cite{lotte2007review, lotte2018review, abiri2019comprehensive} e le mette a disposizione della collettività così da sviluppare una ricca rete internazionale. In particolare, sono tre gli aspetti presi principalmente in considerazione: l’hardware, il software e i paradigmi per l’estrazione dei dati.

\section{Le tecnologie per il tracciamento del segnale elettroencefalografico}
Seppur all’interno della ricerca venga preso in considerazione esclusivamente l’elettroencefalografo, è importante soffermarsi su tutte le strumentazioni disponibili ad oggi. In base al grado di invasività è possibile distinguere due macrocategorie: interfacce neurali invasive e non invasive \cite{peksa2023stateoftheart}. I dispositivi invasivi impiegano microelettrodi intracorticali impiantati direttamente nell’encefalo. Rientrano in questa categoria l’EEG di profondità e l’elettrocorticografia (ECoG). Le tecnologie non invasive si limitano invece a misurare l’attività cerebrale dalla superficie dello scalpo o comunque in maniera indiretta. Sono tecnologie non invasive l’elettroencefalografia (EEG), la magnetoencefalografia (MEG), la risonanza magnetica funzionale e la spettroscopia funzionale nel vicino infrarosso. 

I dispositivi invasivi forniscono un’elevata risoluzione spaziale e temporale, un’eccellente qualità del segnale e la possibilità di effettuare misurazioni altamente localizzate. Per contro, l’invasività della procedura chirurgica necessaria per l’impianto riduce drasticamente la portabilità della metodologia e aumenta notevolmente i costi associati. Il rischio di infezioni o di danni tissutali, nonché il rigetto e il decadimento dei microelettrodi impiantati complicano ulteriormente il quadro.
I dispositivi non invasivi presentano caratteristiche opposte: ridotta risoluzione spaziale, elevata sensibilità ai disturbi, alle interferenze ambientali e agli artefatti muscolari e oculari, bassa localizzazione. Sono caratterizzati, però, da costi relativamente ridotti e facilità d’uso. Nel contesto delle BCI, l’EEG assume un ruolo centrale grazie alla sua elevata portabilità e alla misura che effettua, unica a restituire dati diretti dell’attività neuronale. 

\section{Paradigmi per l’acquisizione dati}
Nei sistemi BCI è necessario sottoporre i soggetti a determinati stimoli e richieste per ricavare informazioni utili e riproducibili. L’insieme di questi attributi definisce un paradigma di acquisizione dei dati. Negli anni sono state proposte differenti tipologie di paradigmi con lo scopo di comprendere il potenziale di differenti canali comunicativi. Ad oggi sono due le metodologie principalmente impiegate: i paradigmi di immaginazione motoria (\emph{Motor Imagery Paradigms}, MIP) e i paradigmi basati su stimolazioni esterne (\emph{External Stimulation Paradigms}, ESP) \cite{abiri2019comprehensive}.

Gli MIP sono basati sull’immaginazione motoria di grandi parti del corpo: durante una sessione con tale paradigma, al soggetto viene richiesto di immaginare un movimento cinestetico senza eseguire realmente tale gesto. Lo scopo è quello di scatenare le \emph{event-related desynchronizations} (ERD), una diminuzione temporanea della potenza di specifiche bande di frequenza, prominente nei ritmi alfa (8-12 Hz) e beta (18-26 Hz), rispetto a un livello di riferimento. Questo livello è generalmente associato allo stato di rilassamento durante il quale si verifica l’evento opposto, le \emph{event-related synchronizations} (ERS). L’obiettivo di tale sistema è quello di tradurre l’intenzione cerebrale di soggetti con limitate capacità motorie in comandi utili. 

In base alla tipologia di classificazione si distinguono due classi: i \emph{sensorimotor rhythms paradigms} sono associati ad una classificazione discreta, all’interno della quale ogni intenzione cerebrale è tradotta in uno specifico comando all’interno di un ristretto insieme di possibilità; gli \emph{imagined body kinematic paradigms} invece fanno uso di un controllo ad elevata dimensionalità con lo scopo di decodificare i parametri cinematici continui del movimento immaginato. Entrambi gli approcci sono caratterizzati da una lenta curva di apprendimento ed una ridotta adattabilità, nonché da un’elevata variabilità fra soggetti: una parte non trascurabile degli utenti non riesce infatti a controllare il sistema, fenomeno noto come \textbf{BCI illiteracy}.

Gli ESP sfruttano invece variazioni dell’attività cerebrale legate a stimolazioni esterne di tipo visivo, uditivo o somatosensoriale. Tali variazioni sono note nella letteratura scientifica come \emph{Event-Related Potentials} (ERP): somministrando lo stesso stimolo un elevato numero di volte ed eseguendo operazioni di \emph{averaging} è possibile visualizzare una risposta cerebrale fisiologica. Non è possibile, in generale, visualizzare la reazione ad un singolo stimolo: per questo motivo i tempi di latenza di questa famiglia di paradigmi sono uno degli ostacoli principali. Anche in questo caso sono comunemente distinte due classi di paradigmi, il \emph{Visual P300 paradigm} e lo \emph{Steady State Visual Evoked Potentials Paradigms} (SSVEP) \cite{maslova2023noninvasive}. Entrambi gli approcci si basano sulla somministrazione di uno stimolo visivo intermittente, presentato in determinati istanti di tempo oppure ad una precisa frequenza. Al soggetto viene richiesto di visualizzare il quadrante luminoso associato al comando che è intenzionato a mettere in atto e l’algoritmo di riconoscimento esegue la classificazione di tale intenzione. Rispetto al precedente paradigma, l’ESP è semplice e immediato e richiede una calibrazione rapida di pochi minuti. Lo sforzo visivo e il grado di attenzione richiesto sono però limiti significativi.

Oltre ai paradigmi presentati, è possibile osservare numerose strategie impiegate in maniera minore. Un approccio molto interessante impiega un meccanismo di feedback sugli errori. Quando l’intenzione del soggetto non corrisponde all’azione eseguita dal sistema BCI si manifesta un \emph{Error-Related Potential} (ErrP) nella corteccia mediale-frontale, con una latenza di 200-700 ms. La ricerca moderna sta implementando questa metodologia principalmente in sistemi autonomi nei quali viene richiesto ai soggetti di supervisionare e correggere eventuali comportamenti errati. Individuare una correlazione tra l’ampiezza del segnale e il corrispettivo errore e riuscire ad utilizzare queste informazioni per ottenere una correzione real-time sono gli aspetti principalmente indagati.

Infine, è possibile utilizzare una qualsiasi combinazione di dispositivi hardware e paradigmi e sviluppare nuove forme di \emph{hybrid paradigms}. L’impiego può essere in serie, con l’output di un paradigma utilizzato in input per il successivo paradigma, oppure in parallelo. Ne sono un esempio gli studi combinati di SSVEP e P300 per l’estrazione contemporanea delle informazioni di direzione e velocità.

\section{Algoritmi per l’estrazione e la selezione di feature}
Come avviene in altre tipologie di pipeline per il machine learning, anche nel caso dei sistemi BCI i metodi di classificazione non si limitano all’uso di un unico classificatore ma necessitano di algoritmi per l’estrazione e la selezione di feature necessari per una rappresentazione compatta e discriminativa dei segnali EEG. 

A seguito delle procedure di pre-processing, che applicano filtri sia temporali che spaziali, è possibile sintetizzare l’informazione contenuta nel segnale elaborato attraverso algoritmi di \emph{feature extraction}. Il risultato di questa procedura è un vettore di feature di dimensionalità $d$. 
È possibile rappresentare i segnali EEG in differenti modi, ognuno dei quali trasporta con sé una frazione dell’intera informazione. 
Le feature più semplici e significative si osservano nel dominio del tempo: media, varianza, deviazione standard, energia del segnale sono parametri immediati e semplici, utili per valutare la qualità del segnale; contengono però un’esigua quantità di informazioni nel contesto BCI e sono poco robuste al rumore. 
Nel dominio delle frequenze è contenuta l’informazione fondamentale della densità spettrale, dalla quale è possibile ricavare la potenza per ogni specifica banda EEG; il motore di questa informazione è l'ormai nota correlazione tra specifiche bande e determinate attività. 
Per non perdere l’informazione temporale è possibile, inoltre, applicare strumenti quali la trasformata di Wavelet o la Short-Time Fourier Transform (STFT) così da estrarre una rappresentazione tempo-frequenziale, altamente indicata per l’EEG per via della sua non stazionarietà, seppur costosa computazionalmente. Infine, molte informazioni sono disponibili all’interno del dominio dello spazio dove si ricercano combinazioni tra canali attraverso il rintracciamento di canali correlati o l’impiego di algoritmi altamente specializzati quali il Common Spatial Patterns (CSP).

Studi moderni suggeriscono un nuovo paradigma secondo il quale è possibile ottimizzare l’informazione attraverso rappresentazioni matriciali o tensoriali, piuttosto che vettoriali \cite{lotte2018review}. Frequentemente questo genere di approccio fornisce feature che sono combinazioni lineari di informazioni contenute in diversi domini e dunque di difficile interpretazione fisica; nonostante ciò, i contributi per il settore delle BCI sono notevoli e dunque non trascurabili.
Prima di eseguire la classificazione è possibile aggiungere un ulteriore step, la \emph{feature selection}: si seleziona un sottoinsieme delle feature precedentemente calcolate con lo scopo di ottenere una rappresentazione più compatta. Le ragioni di questa operazione sono molteplici: ridurre la dimensionalità del problema e, con essa, il rischio di \emph{overfitting}, particolarmente elevato quando il numero di trial disponibili è ridotto rispetto a quello delle feature; eliminare le feature ridondanti o poco informative, che aggiungono soltanto rumore; ridurre il costo computazionale, aspetto rilevante per i sistemi real-time; infine, facilitare l’interpretazione dei risultati, individuando ad esempio i canali o le bande di frequenza più discriminanti \cite{lotte2007review}.
Ultimo elemento della pipeline è il classificatore, che associa a ciascun vettore di feature la classe corrispondente. Negli ultimi anni le reti neurali e, più in generale, gli algoritmi di \emph{deep learning} hanno conosciuto una rapida diffusione anche nel settore delle BCI \cite{hossain2022status}: le reti convoluzionali e ricorrenti sono in grado di apprendere direttamente dal segnale poco elaborato, integrando in un unico modello le fasi di estrazione delle feature e di classificazione. A fronte di questi vantaggi, tali modelli richiedono grandi quantità di dati per l’addestramento e un elevato costo computazionale, e restituiscono risultati di difficile interpretazione. Per questo motivo i classificatori più classici mantengono ancora oggi un ruolo centrale. L’analisi discriminante lineare (LDA) individua l’iperpiano che meglio separa le classi, assumendo per ciascuna una distribuzione gaussiana con matrice di covarianza comune: è semplice, rapida e poco soggetta a overfitting, e la sua variante con contrazione (\emph{shrinkage}) risulta particolarmente adatta quando i campioni disponibili sono pochi rispetto al numero di feature. Le macchine a vettori di supporto (SVM) ricercano invece l’iperpiano che massimizza il margine fra le classi e, attraverso l’impiego di funzioni \emph{kernel}, possono modellare anche confini non lineari. Nei contesti caratterizzati da pochi dati per soggetto, come quello di questa tesi, entrambi offrono prestazioni spesso confrontabili con quelle delle reti neurali a un costo computazionale nettamente inferiore.
